Whether the failures that matter in service will be observable is rarely a criterion when the architecture of an engineered system is chosen; condition monitoring is specified later. This paper treats diagnosability as a property of the triple (architecture, control structure, observation suite) and evaluates it with one scale-free statistic, the signal-to-drift ratio: the relative change of an observable feature under a failure mode of known extent, divided by the relative change the healthy system exhibits on its own. From it follow a minimum detectable extent and a detectable share per alternative, observation suite and failure-mode class, a decision table over the required extent and false-alarm level, and a six-step procedure that co-selects architecture and monitoring suite and tests whether a monitoring design transfers across a product family. The method is demonstrated on three wind-turbine generator topologies under identical winding short-circuits on open benchmark data. The same design-level fault produces different physical severities; where a fault becomes visible depends on the machine--drive combination, and the detectability difference between two machines on one bench comes from their healthy non-stationarity, not the signature; the cheapest suite depends on the alternative and on the requirement, voltage transducers lowering the permanent-magnet machine's detectable extent and adding nothing for the induction machine; and a learned detector transfers across architectures only when its features are referenced to each trial's own preceding healthy behaviour. Results carry bootstrap uncertainty, are limited to one machine per alternative, and are reproducible from the public data with the released code.
